\clearpage
\section{Entrega del Producto - Entregar la aplicación, la biblioteca y su documentación [Pendientes: 1. Acordar las fechas y el medio de entrega. 2. Definir el procedimiento y los responsables del empaquetado y la presentación.]}\label{sec:spmp-11}

Durante la Semi-transición se entregará la versión evaluada de la aplicación de escritorio, su biblioteca de escenarios FDS precalculados y la documentación necesaria para instalar y utilizar la aplicación y revisar sus resultados~\cite{adaptacion2026,std12207}. La entrega incluirá los siete productos del catálogo de la sección \ref{sec:spmp-6-4}:

\begin{enumerate}
    \item \textbf{E1. Requisitos, alcance y trazabilidad:} especificación y relación entre necesidades, funciones y pruebas.
    \item \textbf{E2. Arquitectura y modelos UML:} diseño, modelos y decisiones principales.
    \item \textbf{E3. Paquete reproducible FDS:} entradas preparadas, configuración e informe del piloto.
    \item \textbf{E4. Biblioteca y catálogo:} escenarios precalculados y condiciones asociadas.
    \item \textbf{E5. Aplicación:} código fuente, ejecutable y dependencias identificadas.
    \item \textbf{E6. Pruebas y evaluaciones:} planes, registros, hallazgos y correcciones.
    \item \textbf{E7. Documentación y presentación:} manual de instalación y uso, guía técnica, informe final y demostración.
\end{enumerate}

\textbf{Procedimiento de entrega.} Se realizarán las siguientes actividades:

\begin{enumerate}
    \item \textbf{Consolidar:} reunir E1--E7, identificar sus versiones y verificar su coherencia con el alcance y la configuración de la sección \ref{sec:spmp-12-3}.
    \item \textbf{Verificar:} comprobar la ejecución de la aplicación y los resultados de E6, aplicando los controles de la sección \ref{sec:spmp-12-5} y, una vez acordados, los criterios de aceptación de la sección \ref{sec:spmp-8-3}.
    \item \textbf{Empaquetar:} organizar el código, el ejecutable, los archivos FDS, la biblioteca y los documentos; identificar la estructura y las versiones. El medio de distribución de resultados pesados se definirá según su volumen, sin perder la relación con el catálogo.
    \item \textbf{Comprobar el paquete:} revisar su integridad, la correspondencia del manual y la identificación de escenarios; cuando sea posible, probar instalación y ejecución, registrando hallazgos.
    \item \textbf{Presentar:} demostrar las funciones verificadas, desde la entrada del incendio y la consulta meteorológica hasta la selección de escenarios y su visualización.
    \item \textbf{Registrar:} documentar los productos entregados, versiones, fecha, medio y observaciones recibidas.
\end{enumerate}

La demostración distinguirá el funcionamiento comprobado de la aplicación de la validez científica de las simulaciones. Los escenarios FDS precalculados son apoyo a la interpretación, no predicciones exactas ni sustitutos del criterio del comandante o coordinador. La evaluación con usuarios y la entrega académica conservarán objetivos y evidencias diferenciados.

Las fechas, el medio y los responsables individuales del empaquetado y presentación se acordarán conforme a la calendarización \ref{sec:spmp-9-3-2} y la organización \ref{sec:spmp-8-4-2}. Toda corrección posterior se registrará y comprobará; si afecta requisitos o compromisos, seguirá las secciones \ref{sec:spmp-10-1} y \ref{sec:spmp-6-6}.
